% Plan Sections 5.1, 5.4, 5.5, 5.7 and 5.9.
\section{Method}
\label{sec:method}

\subsection{Definitions}
\label{sec:definitions}
% Plan Section 5.1: unmeasured pairing, the component kinds, synchronisation,
% operation sequence, recall, accuracy, admissible configuration,
% kernel-speedup rule, profitable.
\stub{definitions, plan Section 5.1}

\subsection{The critical-path model}
\label{sec:model}
% Plan Section 5.4: the recursion, then the baselines built from the same
% components (sum, fitted sum, step form, learned, HDBI as a diagnostic,
% constant host cost, nearest measured pairing), how the fitted baselines are
% trained, and the ablations.
\stub{\bunga recursion, plan Section 5.4}

\stub{baselines and how they are fitted}

\subsection{Arms}
\label{sec:arms}
% Plan Section 5.5, including dense-baseline selection (K2). Keep oracle,
% deployable-estimator, CPU-builder and GPU-builder results distinct (rule 2).
\stub{arms, plan Section 5.5}

\subsection{Quality measures}
\label{sec:quality}
% Plan Section 5.7. Non-inferiority at a stated margin only (rule 1).
\stub{accuracy test, recall, H4 power rule, plan Section 5.7}

\subsection{Profitability dispatcher}
\label{sec:dispatcher}
% Plan Section 5.9.
\stub{dispatcher, plan Section 5.9}
